\documentclass{article}
\usepackage[spanish]{babel}
\usepackage[T1]{fontenc}
\usepackage[utf8]{inputenc}
\usepackage{graphicx}
\usepackage{parskip}

\title{Patrones de Dise\~no en mi proyecto: Simulador de combates de Dragon Ball}
\author{Diego}

\date{Mayo 8 2026}

\begin{document}

\maketitle

\tableofcontents
\newpage

\section{Introducci\'on}



Los patrones se dividen en tres grupos principales: los \textbf{creacionales},
que tienen que ver con c\'omo se crean los objetos; los \textbf{estructurales},
que hablan de c\'omo se organizan entre s\'i y los \textbf{de comportamiento},
que definen c\'omo interact\'uan y se comunican.

Voy a explicar los patrones que yo ocupe en mi proyecto, Simulador de peleas de Dragon Ball.

\section{Patrones Creacionales}

Estos patrones se encargan de controlar c\'omo se crean los objetos dentro de
una aplicaci\'on, puede parecer que eso es lo m\'as simple del
mundo (solo haces \texttt{new Objeto()} y listo), pero en proyectos m\'as
grandes eso puede volverse un problema.

\subsection{Singleton}
La idea es que de una clase solo exista una \'unica instancia en todo el programa. Imag\'inate que tienes una
clase que maneja la conexi\'on a la base de datos. No tiene sentido crear
cien conexiones distintas, con una sola alcanza.

En el proyecto que estoy haciendo, el Singleton para el
\texttt{AuthService}, donde guardaba la sesi\'on del usuario. No quer\'ia
que hubiera varios objetos manejando el token al mismo tiempo porque eso
generar\'ia inconsistencias.

La desventaja que le veo es que a veces complica las pruebas, porque el
estado global puede interferir entre tests, pero en general es muy \'util.

\subsection{Factory Method}

Este patr\'on b\'asicamente dice: en lugar de crear un objeto directamente,
delegas esa responsabilidad a un m\'etodo o clase especial que decide qu\'e
tipo de objeto crear.


Lo us\'e de forma informal cuando el backend dec\'idia qu\'e tipo de respuesta
construir dependiendo del resultado de la operaci\'on, no era un Factory
Method puro, pero la idea era similar.

\subsection{Builder}

El Builder sirve cuando tienes objetos que se construyen paso a paso y
que tienen muchos campos opcionales, en vez de tener un constructor con
veinte par\'ametros (que nadie entiende en qu\'e orden van), usas un objeto
intermedio que vas configurando.

En mi proyecto cuando se va construyendo un objeto \texttt{Combate} con
lugar, a\~no, intensidad, tipo, imagen y resultado, podr\'ia usar un Builder
para ir arm\'andolo sin confundirte con los par\'ametros.

Yo no lo implement\'e como tal, pero s\'i lo tuve en mente cuando
dise\~n\'e los formularios de creaci\'on de combates en Flutter, iba llenando
campo por campo antes de mandar todo al backend.

\section{Patrones Estructurales}

Estos patrones se enfocan en c\'omo se relacionan y se organizan las clases
entre s\'i, sirven para que el c\'odigo sea m\'as flexible y f\'acil de mantener
sin tener que cambiar todo cada vez que algo cambia.

\subsection{Repository}

Este fue el que m\'as us\'e en el proyecto, la idea es poner una capa entre
la l\'ogica de la aplicaci\'on y el origen de los datos (que puede ser una
base de datos, una API, lo que sea).

En Flutter tengo una clase \texttt{CombateRepository} que es la \'unica
que sab\'ia c\'omo hablar con el backend, las pantallas no llamaban
directamente al servicio HTTP, sino que le ped\'ian al repositorio, as\'i,
si en alg\'un momento cambiaba la URL o la forma de obtener datos solo
ten\'ia que modificar el repositorio y no tocar todas las pantallas.

Es un patr\'on muy pr\'actico.

\subsection{Adapter}

El Adapter sirve cuando tienes dos partes del sistema que no hablan el
mismo idioma y necesitas un intermediario que traduzca.

En mi proyecto lo vi reflejado en los m\'etodos \texttt{fromJson()} de los
modelos en Dart, el backend devolv\'ia JSON con ciertos nombres de campos,
y el modelo se encargaba de traducirlos al formato que usaba Flutter
internamente.

\subsection{Facade}

El Facade es b\'asicamente crear una clase sencilla que oculte la
complejidad interna, desde afuera ves una interfaz simple, pero por
dentro puede estar haciendo muchas cosas.

MI \texttt{ApiService} del proyecto funcion m\'as o menos as\'i, las
pantallas solo llaman m\'etodos como \texttt{obtenerCombates()} sin
preocuparse por los  HTTP, el manejo del token etc, 
todo eso estaba escondido adentro del ApiService.

Es muy \'util para mantener el c\'odigo de las pantallas limpio y enfocado
solo en mostrar informaci\'on.







\section{Conclusi\'on}

La verdad es que al principio me parec\'ia que los patrones de dise\~no son muy didicles de aplicar en proyectos reales, pero conforme fui entendiendo, haciendo practicas en clas y lo estoy apliacando en mi proyecto
me doy cuenta de que los estaba usando sin saberlo, o que muchas
cosas que hac\'ia por intuici\'on ten\'ian un nombre y una raz\'on de ser.

No creo que haya que memorizar todos los patrones ni intentar usarlos
todos en un mismo proyecto , laa clave es entender el problema que resuelve
cada uno y tenerlos en mente para cuando uno los necesite.



\end{document}


